%!TEX encoding = UTF8
%!TEX root = 0-notes.tex

\chapter{Géométrie euclidienne}
\label{chap:geometrie}

\section{Droites et équations cartésiennes}

\begin{multicols}{2}

Soit $(d)$ une droite \emph{non verticale} qu'on plonge dans un repère orthonormé.
On suppose $(d)$ non verticale pour qu'elle soit la courbe représentative d'une certaine fonction $f$ qu'on souhaite déterminer : \[(d) = \Ccal_f. \]
%Posons en outre $\D = \R$ le domaine de $f$, soit la droite réelle toute entière.

	\begin{center}
	%\includegraphics[page=1, scale=.9]{figures/fig-affines.pdf}
	\includegraphics[page=1, scale=1]{2-figures.pdf}
	\end{center}

%On a supposé $(d)$ non verticale car dans le cas contraire, un certain antécédent $x$ aurait plusieurs images, ce qui n'est pas possible d'après la définition d'une fonction.

Comme on a supposé $(d)$ non verticale, elle coupe nécessairement l'axe des ordonnées.
Posons $(0; b)$ le point d'intersection, où $b\in\R$ est un réel désormais fixé.

%On souhaite trouver la fonction $f$ telle que $(d) = \Ccal_f$, sa courbe représentative.
On rappelle la propriété fondamentale : 
	\begin{align*}
		(x;y) \in \Ccal_f && \iff && y = f(x).
	\end{align*}
Prenons donc un $x \in \R$ non nul, et essayons de déterminer l'ordonnée du point de $(d)$ d'abscisse $x$ afin d'obtenir $f(x)$.

\vspace{30pt}

On se restreint au triangle rectangle dont les sommets sont $(0;b), (x ; b)$, et $(x;y)$ pour appliquer la trigonométrie.

	\begin{center}
	\includegraphics[page=2, scale=1]{2-figures.pdf}
	\end{center}

L'angle $\alpha$ comme défini ci-dessus est constant et mesure la pente de la droite.

	\begin{align*}
		\tan(\alpha) = \frac{|y-b|}{|x|} && \iff && |y-b| = \tan(\alpha) \cdot |x|.
	\end{align*}
Comme la valeur absolue ne fait que supprimer le signe, on en déduit que
	\begin{align*}
		y - b = \pm \tan(\alpha) \cdot x && \iff && y = \pm \tan(\alpha) \cdot x + b,
	\end{align*}
où $\pm$ signifie ``plus ou moins'', c'est-à-dire qu'on ne connait pas \emph{a priori} le signe et que cela n'importe pas pour l'instant.
En posant $a = \pm \tan(\alpha)$, on a trouvé la forme de la fonction $f$ telle que $(d) = \Ccal_f$ :
	\[ f(x) = a\cdot x + b, \]
où $a$ et $b$ sont des nombres réels déterminés entièrement par la droite $(d)$.


\end{multicols}

\newpage

Dans la suite, on supposera que $a, b, c\in\R$ sont des réels fixés tels que $a$ et $b$ ne sont pas tous les deux nuls.

\dfn{équation cartésienne de droite}{
	Une \emphindex{équation cartésienne de droite} est une équation de la forme
		\begin{align}
			ax + by = c. \label{eq:cart}
		\end{align}
	L'ensemble des couples $(x ; y)$ vérifiant l'équation forme une droite :
		\begin{align}
			(d) = \bigset{ (x ; y) \tq ax + by = c }. \label{eq:droite}
		\end{align}
}{dfn:equation-cartesienne}

%\nt{
%	Le fait que $(d)$ soit une droite peut être vu comme une conséquence de la trigonométrie ou encore à partir de la caractérisation comme l'ensemble de translatés d'un point.
%	Nous justifierons que $(d)$ est une droite en le voyant comme la droite passant par un point et orthogonale à un vecteur. 
%}

\exe{}{
	Tracer la droite d'équation $2x + 3y = 0$ et le vecteur $\pvec23$ dans un même repère.
	Que remarque-t-on ?
}{exe:axbyz}{
	TODO
}

\exe{}{
	Tracer la droite d'équation $4x - y = 3$ et le vecteur $\pvec4{-1}$ dans un même repère.
	Que remarque-t-on ?
}{exe:axbyz2}{
	TODO
}


\notations{
	%
	%Considérons le triangle $OAX$ où $O(0;0)$, $A(a;b)$ et $X(x ; y)$.
	Deux vecteurs $u, v\in\R^2$ sont perpendiculaires dès que les droites $(OU)$ et $(OV)$ sont perpendiculaires, où $O(0;0)$, $U(u_1;u_2)$ et $V(v_1 ; v_2)$
}

\thm{}{
	Un couple (x ; y) vérifie $ax+by=0$ si et seulement si les vecteurs $\pvec{x}{y}$ et $\pvec{a}{b}$ sont perpendiculaires.
	\[ ax+by=0 \iff \pvec{x}{y} \perp \pvec{a}{b}. \]
}{thm:ab-orth}

\pf{}{
	Considérons le triangle $OAX$ où $O(0;0)$, $A(a;b)$ et $X(x ; y)$.
%	Les vecteurs $\pvec{x}{y}$ et $\pvec{a}{b}$ sont perpendiculaires si et seulement si les droites $(OA)$ et $(OX)$ le sont.
	
	D'après le théorème de Pythagore (et sa réciproque), nous avons
		\[ (OA) \perp (OX) \iff OA^2 + OX^2 = AB^2. \]
	Or $OA^2 = a^2 + b^2$, $OB^2 = x^2 + y^2$, et $AB^2 = (a-x)^2 + (b-y)^2$.
	L'égalité de Pythagore devient donc
		\begin{align*}
			(OA) \perp (OX) &\iff a^2 + b^2 + x^2 + y^2 = (a-x)^2 + (b-y)^2, \\
				&\iff a^2 + b^2 + x^2 + y^2 = a^2 + x^2 - 2ax + b^2 + y^2 - 2by, \\
				&\iff 0 = -2(ax + by), \\
				&\iff ax + by = 0.
		\end{align*}
}

%\exe{difficulty=1}{
%	Considérons la droite $(d) : 3x-2y = 5$ du plan cartésien.
%	Trouver l'équation d'une droite $(d')$ perpendiculaire à $(d)$.
%}{exe:droite-perp}{
%	todo
%}

\exe{}{
	Tracer les droites d'équations $2x - y = 0$, $2x - y = -2$, et $2x - y = 3$ dans un même repère.
	Que remarque-t-on ?
}{exe:parallélisme}{
	TODO
}

\exe{, difficulty = 1}{
	Montrer que les droites d'équation $ax+by =0$ et $ax+by=c$ sont parallèles, quel que soit $c\in\R$.	
}{exe:parallélisme2}{
	Si $c=0$, alors les droites sont confondues et donc parallèles.
	
	Sinon, l'intersection des deux droites est vide et elles sont parallèles.
%	Si $b=0$, les droites sont verticales et parallèles.
%	
%	Sinon, leur équations réduites sont $y = \frac{-a}{b}x$ et $y = \frac{-a}{b}x + \frac{c}b$, de même coefficient directeur.
%	Les droites sont donc parallèles dans ce cas aussi.
}

\thm{}{
	Considérons un point $P(x_P;y_P)$ vérifiant $ax_P + bx_P = c$.
	
	Alors la droite d'équation $ax+by = c$ est la droite perpendiculaire au vecteur $\pvec{a}{b}$ et passant par $P$.
}{thm:shift}

\pf{}{
	La droite passe bien sûr par $P$ par construction de celle-ci.
	
	De plus, elle est parallèle à $ax+by=0$ d'après l'exercice \ref{exe:parallélisme2}, et donc perpendiculaire à $\pvec{a}{b}$ d'après le théorème \ref{thm:ab-orth}.
	% mieux argumenter pour que ça soit généralisable plus facilement en 3d ?
}

\exe{difficulty=1}{
	Considérons la droite $(OA)$ où $A(a;b)$.
	
	Montrer que la fonction $f(x,y) = ax+by$ sur $(OA)$ est croissante dans la direction $\vec{OA}$, nulle en $O$, et décroissante dans la direction $\vec{AO}$.
	En particulier, pour n'importe quel $c\in\R$, il existe un point $P(x_P;y_P)$ de cette droite vérifiant
		\[ ax_P + by_P = c. \]
}{exe:allc}{
	On rappelle que la droite $(OA)$ est la droite passant par $O$ et supportée par le vecteur $\vec{OA} = \pvec{a}b$.
	Graphiquement, pour obtenir $(OA)$, on translate l'origine $O$ par tous les multiples du vecteur $\vec{OA}$.
	C'est donc l'ensemble des translatés $O + k\pvec{a}b = (ka ; kb)$, où $k$ parcourt $\R$.

	Le point $(x_P ; y_P) = (ka ; kb)$, parcourant $(OA)$ et passant par l'origine, vérifie
		\[ ax_P + by_P = k(a^2 + b^2). \]	
	Comme $a$ et $b$ ne sont pas tous les deux nuls, $a^2 + b^2 >0$. 
	L'expression $ax_P + by_P$ est ainsi une fonction affine croissante en $k$ : quand $k$ grandit, $ax_P + by_P$ aussi.
	
	Pour tout $c\in\R$, choisir $k = \frac{c}{a^2 + b^2}$ permet bien d'obtenir $ax_P + by_P = c$.
}

\exe{}{
	Tracer la droite $(OA)$ où $A(2, -1)$. Calculer les coordonées des points $(x;y)$ de $(OA)$ d'abscisse $-2, -1, 0, 1,$ et 2, et les placer sur la droite.
	Pour chacun des points $(x;y)$ obtenus, calculer $2x -y$.
}{exe:uscalaru}{
	TODO
}

\section{Géométrie euclidienne dans $\R^n$}

\dfn{espace $\R^n$ et opérations}{
	Soit $n\in\N^*$ un entier naturel.
	L'ensemble des \emphindex{vecteurs} $v$ à $n$ coordonnées réelles, dénotés
		\begin{align*}
			 v = \begin{pmatrix} v_1 \\ \vdots \\ v_n \end{pmatrix} && v_i \in \R \quad \forall i=1, \dots, n,
		\end{align*}
	est noté $\R^n$ et lu « R $n$ ».
	
	Il est possible d'ajouter des vecteurs \underline{de même dimension} ensemble, et de multiplier un vecteur avec un réel alors appelé \emphindex{scalaire}.
	\begin{multicols}{2}
	\begin{enumerate}
		\item
		$u+v = \begin{pmatrix} u_1 \\ \vdots \\ u_n \end{pmatrix} + \begin{pmatrix} v_1 \\ \vdots \\ v_n \end{pmatrix} = \begin{pmatrix} u_1+v_1 \\ \vdots \\ u_n+v_n \end{pmatrix}$.
		\item
		$\alpha u = \alpha \begin{pmatrix} u_1 \\ \vdots \\ u_n \end{pmatrix} = \begin{pmatrix} \alpha u_1 \\ \vdots \\ \alpha u_n \end{pmatrix}$
	\end{enumerate}
	\end{multicols}
}{dfn:Rn}

\dfn{indépendance linéaire}{
	De façon analogue à la définition \ref{dfn:indep-lin-poly}, des vecteurs $v_1 ; v_2 ; \dots ; v_k \in \R^n$ sont linéairement indépendants dès que
	\begin{align*}
		\alpha_1 v_1 + \alpha_2 v_2 + \dots + \alpha_k v_k = 0 && \implies && \alpha_1=\alpha_2=\dots=\alpha_k = 0.
	\end{align*}
}{dfn:indep-lin-vec}

\prop{}{
	De façon analogue à la proposition \ref{prop:indep-unicite-poly}, si $v$ est combinaison linéaire de vecteurs linéairement indépendants, alors cette combinaison est unique.
}{prop:indep-unicite-vec}

\nomen{
	Deux vecteurs sont \emphindex{colinéaires} s'ils sont multiples l'un de l'autre.
}

\notation{
	Dans la suite, nous noterons les coordonnées d'un vecteur $u\in\R^n$ par
		\[ u = \begin{pmatrix}u_1 \\ u_2 \\ \vdots \\ u_n \end{pmatrix}. \]
}

\ex{}{
	\begin{multicols}{2}
	\begin{enumerate}
		\item
		$\pvec12 \in \R^2$
		\item
		$\R^1 = \R$
		\item
		$-2\begin{pmatrix}2 \\ -\pi \\ \frac12 \end{pmatrix} = \begin{pmatrix}-4 \\ 2\pi \\ -1 \end{pmatrix} \in \R^3$
		\item
		$3\begin{pmatrix}2\\-4\\3\end{pmatrix} + \begin{pmatrix}2\\4\\-1\end{pmatrix} = \begin{pmatrix}3\times2+2\\3\times(-4)+4\\3\times3+(-1)\end{pmatrix} =\begin{pmatrix}8\\-8\\8\end{pmatrix}$
		\item
		$1 = \begin{pmatrix} 1 \end{pmatrix} \in \R^1$
	\end{enumerate}
	\end{multicols}
}{ex:Rn}

\subsection{Produit scalaire}

\dfn{produit scalaire}{
	Soient $u$ et $v$
	deux vecteurs de $\R^n$.
	Le \emphindex{produit scalaire} entre $u$ et $v$ est donné par
		\[ \japb{u, v} = \sum_{i=1}^n u_iv_i =  u_1 v_1 + u_2v_2 + \dots + u_nv_n. \]
}{dfn:produit-scalaire}

\notation{
	Le produit scalaire entre $u$ et $v$ est lu « $u$ scalaire $v$ » et admet plusieurs notations :
	%\begin{multicols}{2}
	\begin{enumerate}
		\item
		$u \cdot v$, d'où le \emphindex{dot product} ;
		\item
		$\scal{u, v}$, utilisée pour les formes bilinéaires ;
		\item
		$u^T v$, produit matriciel ;
		\item
		$u_iv_i$, convention de sommation d'Einstein.
	\end{enumerate}
	%\end{multicols}
}


\thm{bilinéarité du produit scalaire}{
	Le \emphindex{produit scalaire} $\scal{. , .}$ est une \emphindex{forme bilinéaire symétrique définie positive}.
	Pour $x, y, z\in\R^n$ et $\alpha, \beta \in \R$, la forme est 
	\begin{enumerate}
		\item linéaire à gauche : $\japb{\alpha x + \beta y, z} = \alpha\japb{x, z} + \beta \japb{y, z}$ ;
		\item linéaire à droite : $\japb{x, \alpha y + \beta z} = \alpha\japb{x, y} + \beta \japb{x, z}$ ; 
		\item symétrique : $\japb{x, y} = \japb{y, x}$ ;
		\item définie : $\japb{x, x} = 0$ si et seulement si $x=0$ ;
		\item positive : $\japb{x, x} \geq 0$.
	\end{enumerate}
}{thm:bilinearite-scalaire}

\pf{}{
	La symétrie est claire, et il suffit donc de montrer la linéarité à gauche.
	Pour $x, y, z\in\R^n$ et $\alpha, \beta \in \R$, nous avons
		\[  \japb{\alpha x + \beta y, z} = (\alpha x + \beta y)_i z_i = (\alpha x_i + \beta y_i) z_i = \alpha x_i z_i + \beta y_i z_i =  \alpha\japb{x, z} + \beta \japb{y, z}. \]
}


\dfn{orthogonalité}{
	Deux vecteurs de $\R^n$ sont dits \emphindex{vecteurs orthogonaux} dès que $\scal{u, v} = 0$.
}{dfn:orthogonalite}


\exe{}{
	Pour chaque paire de vecteurs suivante, sont-ils orthogonaux ?
	\begin{multicols}{3}
	\begin{enumerate}
		\item $\pvec12$ et $\pvec{-1}{-2}$
		\item $\pvec12$ et $\pvec{-2}{1}$
		\item $\pvec{-6}{-3}$ et $\pvec{1}{-2}$
		\item $\pvec{-6}{-3}$ et $\pvec{-1}{2}$
		\item $\begin{pmatrix}2 \\ 4 \\ -3 \end{pmatrix}$ et $\begin{pmatrix}5 \\ -1 \\ -2 \end{pmatrix}$
		\item $\begin{pmatrix}2 \\ 4 \\ -3 \end{pmatrix}$ et $\begin{pmatrix}-5 \\ 1 \\ 2 \end{pmatrix}$
	\end{enumerate}
	\end{multicols}
}{exe:is-orth}{
	todo
}

\exe{difficulty=1}{
	Soit $a\in\R$. Donner un vecteur orthogonal à $\pvec1a$.
}{exe:vec-orth}{
	todo
}

\exe{difficulty=1}{
	Si $u, v\in\R^2$ sont orthogonaux, montrer que $u$ et $-v$ le sont aussi puis que $u$ et $\lambda v$ le sont aussi pour tout $\lambda\in\R$.
	Interpréter géométriquement.
}{exe:vec-orth-lin}{
	todo
}


\subsection{Norme}



\dfn{norme}{
	%Soit $u = \begin{pmatrix}u_1 \\ u_2 \\ \vdots \\ u_n \end{pmatrix}$
	Soit $u$
	un vecteur de $\R^n$.
	La \emphindex{norme} de $u$ est donnée par
		\[ \norm{u}^2 = u_1^2 + u_2^2 + \dots + u_n^2. \]
}{dfn:norme}

\nt{
	Le produit scalaire d'un vecteur $u$ avec lui-même égale le carré de sa norme.
		\[ \scal{u, u} = \norm{u}^2. \]
}

\prop{}{
	La longueur du segment $[AB]$ est donné par $\left\Vert\vec{AB}\right\Vert$.
}{prop:norm-l2}

\pf{}{
	Théorème de Pythagore.
}

\exe{}{
	Calculer la distance entre les points $P(3;4)$ et $Q(-2;2)$ dans un repère orthonormé.
}{exe:dist-2d}{
	todo
}

\exe{}{
	Soient $u, v$ deux vecteurs de $\R^n$. Montrer l'identité remarquable
		\[ \norm{u+v}^2 = \norm{u}^2 + \norm{v}^2 + 2\scal{u, v}. \]
}{exe:id-rem-scal}{
	todo
}

\notation{
	Dans ce qui suit, nous pourrons voir un point comme un vecteur, et lui appliquer toute la structure des vecteurs (somme, produit par un scalaire, produit scalaire, etc...).
	
	Ainsi, nous pouvons écrire, pour deux points $A, B\in\R^n$, que la longueur du segment $[AB]$ est donnée par
		\[ AB = \norm{\vec{AB}} = \norm{B - A} = \norm{A - B}. \]
}

% inégalités triangulaire et Cauchy-Schwarz sont Algèbre IV
%\thm{inégalité triangulaire}{
%	La norme euclidienne vérifie l'\emphindex{inégalité triangulaire} :
%	pour tout $u, v \in\R^n$,
%		\[ \norm{u + v} \leq \norm{u} + \norm{v}. \]
%}{thm:ineg-tri}


\dfn{}{
	Une famille $\{ v_1 ; \dots ; v_k \} \subseteq \R^n$ est \emphindex{orthonormée} ou \emphindex{orthonormale} dès que
	\begin{enumerate}
		\item
		elle est \emphindex{orthogonale} :
			\[ \japb{v_i, v_j} = 0 \qquad \forall i \neq j ; \]
		\item
		les vecteurs sont \emphindex{unitaires} :
			\[ \norm{u_i} = 1 \qquad \forall i. \]
	\end{enumerate}
}{dfn:orthonormee}

\notation{
	Une famille orthonormée vérifie donc
		\[ \japb{v_i, v_j} = \delta_{ij}, \]
	où $\delta_{ij} = \begin{cases*} 1 & si $i=j$, \\ 0 & si $i\neq j$. \end{cases*}$ est le \emphindex{symbole de Kronecker}.
}

\nt{
	Il est facile de passe d'une famille orthogonale composée de vecteurs non nuls à une famille orthonormée : il suffit de \emphindex{normaliser} les vecteurs, c'est-à-dire les diviser par leur norme.
}

\ex{}{
	Dans $\R^n$, la famille composée des $n$ vecteurs suivants est orthonormée.
		\[ \left\{ \begin{pmatrix} 1 \\ 0 \\ 0 \\ \vdots \\ 0 \end{pmatrix} ; \begin{pmatrix} 0 \\ 1 \\ 0 \\ \vdots \\ 0 \end{pmatrix} ; \begin{pmatrix} 0 \\ 0 \\ 1 \\ \vdots \\ 0 \end{pmatrix} ; \dots ; \begin{pmatrix} 0 \\ 0 \\ 0 \\ \vdots \\ 1 \end{pmatrix} \right\} \] 
	Le $i$-ème vecteur ($i=1, 2, \dots, n$) est construit, rempli avec des $0$, et avec un unique $1$ en $i$-ème position.
	Ces vecteurs s'appellent des \emphindex{vecteurs canoniques}. La famille s'appelle la \emphindex{base canonique}.
}{ex:base-canonique}

\thm{de Pythagore}{
	Une famille $\{ v_1 ; \dots ; v_k \} \subseteq \R^n$ orthogonale de vecteurs non nuls est linéairement indépendante sur $\R$.
	
	De plus, pour tous $\alpha_1, \dots, \alpha_k \in \R$ des nombres réels,
		\[ \norm{ \alpha_1 v_1 + \dots + \alpha_k v_k}^2 = \alpha_1^2\norm{v_1}^2 + \dots + \alpha_k^2\norm{v_k}^2 . \]
}{thm:Pythagore}

\pf{}{
	Toute combinaison linéaire nulle prend la forme
		\[  \alpha_1 v_1 + \dots + \alpha_k v_k = 0 \]
	pour certains $\alpha_i\in\R$ réels.
	
	Le produit scalaire de cette combinaison avec un vecteur $v_i$ est égal à
		\[ \japb{  \alpha_1 v_1 + \dots + \alpha_k v_k , v_i } = \japb{0, v_i }. \]
	Or, $\japb{0, v_i} = 0$ et, par linéarité du produit scalaire,
		\[ \japb{  \alpha_1 v_1 + \dots + \alpha_k v_k , v_i } = \alpha_1 \japb{v_1, v_i} + \dots + \alpha_k \japb{v_k, v_i} = \alpha_i, \]
	car tous les produits scalaires sont nuls sauf $\japb{v_i, v_i}$ qui est non nul car $v_i$ est supposé non nul.
	Par conséquent, $\alpha_i = 0$ pour tout $i=1, \dots, k$.
	L'unique combinaison linéaire nulle est la combinaison triviale et les vecteurs sont linéairements indépendants.
	
	Montrons désormais l'identité de Pythagore.
	Comme $\norm{u}^2 = \japb{u, u}$, nous avons, pour tous $\alpha_1, \dots, \alpha_k \in \R$ des nombres réels,
		\[ \norm{ \alpha_1 v_1 + \dots + \alpha_k v_k}^2 = \japb{\alpha_1 v_1 + \dots + \alpha_k v_k, \alpha_1 v_1 + \dots + \alpha_k v_k}. \]
	Par bilinéarité du produit scalaire et orthogonalité de la famille, cette quantité est égale à
		\[ \japb{ \sum_{i=1}^k \alpha_iv_i, \sum_{j=1}^k \alpha_jv_j} = \sum_{i, j = 1}^{k} \alpha_i \alpha_j \japb{v_i, v_j} = \sum_{i=1}^k \alpha_i^2 \norm{v_i}^2. \]
}

\subsection{Application : projeté orthogonal}

Le but de cette section est de trouver un point $Q$ appartenant à une droite ou un plan et le plus proche d'un point $P$ donné.
L'unicité d'un tel point n'est pas claire \apriori, mais nous verrons qu'il est bien unique et que c'est un projeté orthogonal.


\prop{}{
	Soit $P\in\R^n$ un point, $v\in\R^n$ un vecteur non nul, et $\R v = \{ k v \tq k \in \R \}$ la droite engendrée par $v$.
	Soit $Q$ le point le plus proche de $P$ appartenant à la droite. Autrement dit,
		\[ Q = \arg\min_{k\in\R} \norm{ P - kv }. \]
	Alors
		\[ Q = \frac{\japb{P, v}}{\japb{v, v}} v. \]
}{prop:point-sur-droite}

\pf{}{
	Cherchons le nombre $\lambda\in\R$ tel que $P-\lambda v$ soit perpendiculaire à $v$.
	\begin{align*}
		\japb{P - \lambda v, v} &= 0 \\
		\japb{P, v} - \lambda \japb{v, v} &= 0 \\
		\lambda &= \frac{\japb{P, v}}{\japb{v, v}}
	\end{align*}
	Posons $w=P - \lambda v$.
	Soit $Q = \mu v$, où $\mu\in\R$ est un nombre, un point quelconque de la droite $\R v$.
	
	Par construction, $w$ et $v$ sont orthogonaux, donc le théorème de Pythagore \ref{thm:Pythagore} nous donne 
	\begin{align*}
		\norm{P-Q}^2 &= \norm{w + \lambda v - \mu v}^2 \\
				&= \norm{w + (\lambda-\mu)v}^2 \\
				&= \norm{w}^2 + (\lambda-\mu)^2 \norm{v}^2
	\end{align*}
	Quelque soit $\mu\in\R$, le carré  $(\lambda-\mu)^2$ est positif, et est nul uniquement en $\mu = \lambda$.
	
	La distance $\norm{P-Q}^2$ atteint donc son unique minimum en $\mu=\lambda$, et $Q^* = \lambda v$ est bien le point le plus proche de $P$ appartenant à la droite $\R v$.
}

\nomen{
	Le coefficient $\frac{\japb{P, v}}{\japb{v, v}}$ s'appelle \emphindex{coefficient de Fourier}.
	Nous le retrouverons lors de l'étude des séries de Fourier en analyse.
}

\nt{
	Une droite qui ne passe pas par l'origine prend la forme $R + \R v$ avec $R$ un point et $v$ un vecteur.
	Trouver un point $Q$ de la droite le plus proche de $P$.
	Le point $Q$ minimise donc $\norm{ P - (R+ k v)} = \norm{ (P-R) - kv}$ pour $k\in\R$.
	Ceci est entièrement équivalent à trouver le point le plus proche de $P-R$ appartenant à la droite $\R v$.
	La proposition \ref{prop:point-sur-droite} s'applique donc.
	L'interprétation géométrique est la suivante : en décalant tout l'espace de $R$, la droite passe par l'origine et le point $P$ devient $P-R$.
}

\prop{}{
	Soit $P\in\R^n$ un point, $v, w\in\R^n$ deux vecteurs orthogonaux non nuls, et $\R v + \R w = \{ k v + k' w \tq k, k' \in \R \}$ le plan engendré par $v$ et $w$.
	Soit $Q$ le point le plus proche de $P$ appartenant à la droite. Autrement dit,
		\[ Q = \arg\min_{P' \in \R v + \R w} \norm{ P - P' }. \]
	Alors
		\[ Q =\frac{\japb{P, v}}{\japb{v, v}} v + \frac{\japb{P, w}}{\japb{w, w}} w. \]
}{prop:point-sur-plan}

\pf{}{
	Soit $Q =\frac{\japb{P, v}}{\japb{v, v}} v + \frac{\japb{P, w}}{\japb{w, w}} w$.
	Montrons que $P-Q$ est orthogonal à $v$ et à $w$.
	D'abord, comme $\japb{w, v} = 0$ par orthogonalité supposée de $v$ et $w$,
	\begin{align*}
		 \japb{Q, v} &= \japb{ \frac{\japb{P, v}}{\japb{v, v}} v + \frac{\japb{P, w}}{\japb{w, w}} w, v} \\
		 		&= \frac{\japb{P, v}}{\japb{v, v}} \japb{v, v} + \frac{\japb{P, w}}{\japb{w, w}} \japb{w, v} \\
		 		&= \japb{P, v}
	\end{align*}
	il suit donc que $\japb{P-Q, v} = \japb{P, v} - \japb{Q, v} = 0$.
	Ensuite, $\japb{P-Q, w} = 0$ de façon similaire.
	
	Posons désormais $\alpha_v = \frac{\japb{P, v}}{\japb{v, v}}$ et $\alpha_w = \frac{\japb{P, w}}{\japb{w, w}}$ pour simplifier les notations.
	Nous avons toujours $Q = \alpha_v v + \alpha_w w$.
	
	Alors, pour tout $kv + k'w$ appartenant au plan $\R v + \R w$, où $k, k'\in\R$ sont des nombres, nous avons
	\begin{align*}
		\norm{P - kv - k'w}^2 &= \norm{P - Q + (\alpha_v - k)v + (\alpha_w - k')w}^2 \\
				&= \norm{P-Q}^2 + (\alpha_v - k)^2 \norm{v}^2 + (\alpha_w - k')^2 \norm{w}^2
	\end{align*}
	d'après le théorème de Pythagore \ref{thm:Pythagore}.
	
	Cette quantité est minimale uniquement lorsque $k=\alpha_v$ et $k'=\alpha_w$, en utilisant la positivité du carré $x^2 \geq 0$ pour $x\in\R$, égal à zéro uniquement en $x=0$.
}

\exe{}{
	Soient $u = \begin{pmatrix}1\\3\\-4\end{pmatrix}$ et $v = \begin{pmatrix} 1 \\ -1 \\ 0\end{pmatrix}$.
	Quel $\lambda\in\R$ choisir pour que les vecteurs $v$ et $u-\lambda v$ soient orthogonaux ?
	
	En déduire le point de plus proche de $P(1;3;-4)$ appartenant à la droite $\R v$.
}{exe:Fourier-coeff}{

}

\exe{}{
	Trouver le point le plus proche de $P(2;-1;0)$ appartenant à $(2; 1; -1) + \R \begin{pmatrix}2\\0\\-1\end{pmatrix}$.
}{exe:proj-orth-R3}{

}

\exe{}{
	Soient $u = \begin{pmatrix}34\\17\\-17\end{pmatrix}$, $v = \begin{pmatrix} 0 \\ -1 \\ -1\end{pmatrix}$ et  $w = \begin{pmatrix} 2 \\ -2 \\ 1\end{pmatrix}$.
	Quels $\lambda, \mu\in\R$ choisir pour que fois les vecteurs $v$ et $u-\lambda v-\mu w$ soient orthogonaux et les vecteurs $w$ et $u-\lambda v -\mu w$ soient orthogonaux ?
	Poser un système de deux équations à deux inconnues puis le résoudre.
	
	En déduire le point de plus proche de $P(34;17;-17)$ appartenant au plan $\R v + \R w$.
}{exe:proj-orth-plan}{

}

\exe{difficulty=1}{
	Montrer que la famille $\mathcal{F}$ suivante est orthonormée.
		\[ \mathcal{F} = \left\{ \frac1{\sqrt6}\begin{pmatrix} 2 \\ 2 \\ 1 \end{pmatrix} ; \frac1{\sqrt6}\begin{pmatrix} -1 \\ 2 \\ -2 \end{pmatrix} ; \frac1{\sqrt6}\begin{pmatrix} -2 \\ 1 \\ 2 \end{pmatrix} \right\} \]
	Exprimer le point
		\[ P(5 ;  5 ; -2 ) \]
	comme combinaison linéaire des vecteurs de $\mathcal{F}$.
	
	\emph{Indication : ignorer les $\frac1{\sqrt6}$ pour simplifier les calculs et les rajouter à la fin.}
		
	En déduire $Q$, le point le plus proche de $P$ appartenant au plan $\R\begin{pmatrix} 2 \\ 2 \\ 1 \end{pmatrix} + \R\begin{pmatrix} -2 \\ 1 \\ 2 \end{pmatrix}$.
}{exe:proj-orthonormal}{

}

\section{Systèmes : la dualité entre intersections et combinaisons linéaires}

\nt{
	Les chapitres précédents nous ont amenés à résoudre certains systèmes, notamment pour vérifier ou invalider une indépendance linéaire.
	Lorsque nous vérifiions l'indépendance linéaire des trois vecteurs $\begin{pmatrix} 1 \\ 2\end{pmatrix}, \begin{pmatrix}3\\4\end{pmatrix}, et \begin{pmatrix}5\\6\end{pmatrix}$, nous résolvons le système
		\[ \systeme[sort={\alpha,\beta,\gamma}]{\alpha + 3\beta + 5\gamma = 0, \\ 2\alpha + 4\beta + 6\gamma = 0}. \]
	La \emphindex{dualité} est le fait de voir le système, non pas en colonnes, mais en lignes.
	Changeons les noms des variables. Le système devient
		\[ \systeme{x + 3y + 5z = 0, \\ 2x + 4y + 6z = 0}. \]
	Voyons désormais la première équation comme celle du plan de $\R^3$ orthogonal à $\begin{pmatrix}1\\3\\5\end{pmatrix}$, et la deuxième équation comme celle du plan de $\R^3$ orthogonal à $\begin{pmatrix}2\\4\\6\end{pmatrix}$.
	Géométriquement, l'intersection de deux plans aboutit à une droite, et le caractère infini du nombre de solutions du système devient clair (et donc l'existence d'une solution non triviale).
	En particulier, trois vecteurs de $\R^2$ ne peuvent pas être linéairement indépendants.
}

Nous verrons, dans un chapitre futur, la proposition suivante.

\prop{}{
	Si  la famille $\{ v_1 ; \dots ; v_k \} \subseteq \R^n$ est linéairement indépendante sur $\R$, alors $k\leq n$.
}{prop:k-leq-n}

\nomen{
	Une \emphindex{heuristique} est une règle pragmatique qui fonctionne la plupart du temps.
	Elle donne, par exemple, une stratégie de recherche ou un aspect de solution, sans être toujours optimale, ou vraie.
}

\heur{
	En voyant un système d'équations comme certaines contraintes imposées aux variables libres, on obtient l'heuristique
		\begin{center}
			( degrés de liberté ) = ( nombre de variables ) -- ( nombre de contraintes )
		\end{center}
	Le degré de liberté donne la forme des solutions : aucun degré de liberté correspond à une unique solution ; un degré de liberté correspond à une droite ; deux degrés à un plan ; etc...
}

\ex{}{
	L'équation $y = x + 1$, à deux variables et une contrainte, a un degré de liberté : c'est une droite.
	Géométriquement, une équation cartésienne est $-x+y = 1$, et la droite est donc perpendiculaire au vecteur $\pvec{-1}1$.
	
	Considérons désormais le système à deux variables et deux contraintes suivant.
		\[ \begin{cases*} a = 2b + 3, \\ 2a + b = -1. \end{cases*} \]
	Géométriquement, ceci correspond à l'intersection de deux droites d'équations cartésiennes $x-2y=3$ et $2x+y=-1$.
	Heuristiquement, ce système n'a aucun degré de liberté : il n'y a \emph{probablement} qu'une seule solution.
	Résoudre le système permet de nous en assurer.
	
	Considérons maintenant le système à trois variables et deux contraintes suivant.
		\[ \systeme{3x-2y+7z = 1, \\ x + y + 2z = 0.} \]
	Géométriquement, ceci correspond à l'intersection de deux plans, le premier orthogonal au vecteur $\begin{pmatrix}3\\-2\\1\end{pmatrix}$, et le deuxième orthogonal au vecteur $\begin{pmatrix}1\\1\\2\end{pmatrix}$.
	Nous nous attendons donc à obtenir un solution en forme de droite.
	Heuristiquement, ce système a un degré de liberté (3 variables moins 2 contraintes).
	
	Vérifions que c'est bien le cas algébriquement.
	\begin{align*}
		\systeme{3x-2y+7z = 1, \\ x + y + 2z = 0.} &&\iff && \systeme{x + y + 2z = 0, \\ 3x-2y+7z = 1.}
		&& \iff && \systeme{x + y + 2z = 0, \\ -5y+z = 1.} \\ && \iff && \systeme{x +11y = -2, \\ -5y+z =  1.}  
		&& \iff && \systeme{x = -11y-2, \\ z = 5y+1.}
	\end{align*}
	L'ensemble des solutions est donc donné par
	
	\begin{align*}
		\left\{ \begin{pmatrix} -11y-2 \\ y \\ 5y+1 \end{pmatrix} \tq y\in\R \right\} &= \left\{ \begin{pmatrix} -2 \\ 0 \\ 1 \end{pmatrix} + \begin{pmatrix} -11 \\ 1 \\ 5 \end{pmatrix}y \tq y\in\R \right\}
		\\	&= \begin{pmatrix}-2\\0\\1\end{pmatrix} + \begin{pmatrix}-11\\1\\5\end{pmatrix} \R.
	\end{align*}
	C'est donc bien une droite, celle engendrée par le vecteur $\begin{pmatrix}-11\\1\\5\end{pmatrix}$, et passant par le point $(-2;0;1)$.
	
	{Vérification : les points en $y=0$ et $y=1$ sont $(x ; y ; z) = (-2 ; 0 ; 1)$ et $(x ; y ; z) = (-13 ; 1 ; 6)$ et sont bien solutions du système.}
}{ex:heuristique}

\ex{}{
	L'heuristique ne fonctionne pas tout le temps. Un exemple clair est le doublement d'équations.
	Deux équations identiques dans un système n'en valent qu'une.
	Par exemple, le système
		\[ \systeme{2x+3y = 1, \\ 4x+6y=2. } \]
	correspond à une droite, car les deux équations sont équivalentes.
	Géométriquement, les vecteurs orthogonaux $\pvec23$ et $\pvec46$ sont colinéaires : ils sont multiples l'un de l'autre.
	%\emphindex{colinéaires} : ils sont multiples l'un de l'autre.
	Remarquons d'ailleurs que, en voyant le système par colonnes par dualité, les vecteurs $\pvec24$ et $\pvec36$ sont également colinéaires.
	
	Généralement, dans un système, si une équation peut être obtenue à l'aide d'une combinaison linéaire des autres, elle est redondante.
}{ex:heuristique2}

%Heuristique degrés liberté = variable - équations.
%
% dualité : v_1, .. v_n indep lin
% 

% bien mais il faut la dimension. dim n => tout R^n
%\exe{}{
%	Soient $v_1, \dots v_n \in\R^n$.
%	Montrer que le système $\alpha_1 v_1 + \dots + \alpha_n v_n = 0$ admet une unique solution si et seulement si le système $\alpha_1 v_1 + \dots + \alpha_n v_n = v$ admet une unique solution pour tout $v\in\R^n$.
%}{exe:unique-sol-homogene}{
%	Les vecteurs sont linéairement indépendants si et seulement si les combinaisons linéaires sont uniques.
%}

\section*{Exercices supplémentaires}
\addcontentsline{toc}{section}{Exercices supplémentaires}


\exe{, difficulty=0}{
	Pour chaque système d'équations d'inconnues $x, y \in \R$, donner un système équivalent tel que les coefficients multipliant $x$ soient opposés (c'est-à-dire leur somme soit nulle).
	\setlength{\columnsep}{1cm}
	\begin{multicols}{2}
	\begin{enumerate}[label=\roman*), itemsep=20pt]
		\item $\systeme{x + 3y = 3, \\  -2x + 2y = 2.}$ % \iff \systeme{\,,\,}$
		\item $\systeme{x - 7y = 1, \\  -3x + 5y = 13.}$ % \iff \systeme{\,,\,}$
		\item $\systeme{2x - y = 3, \\  -x + 2y = -5.}$ % \iff \systeme{\,,\,}$
		\item $\systeme{4x  - y = -4, \\   x + 5y = 2.}$ % \iff  \systeme{\,,\,}$
		\item $\systeme{-2x - 2y = -5, \\  3x + 7y = 10.}$ % \iff\systeme{\,,\,}$
		\item $\systeme{7x - y = 10, \\  5x + 2y = 1.}$ % \iff\systeme{\,,\,}$
	\end{enumerate}
	\end{multicols}
}{exe:2}{
	\setlength{\columnsep}{1cm}
	\begin{multicols}{2}
	\begin{enumerate}[label=\roman*), itemsep=20pt]
		\item $\systeme{x + 3y = 3, \\  -2x + 2y = 2.}$ % \iff \systeme{2x+6y=6, \\ -2x+2y=2.}$
		\item $\systeme{x - 7y = 1, \\  -3x + 5y = 13.}$ % \iff \systeme{3x - 21y = 3, \\  -3x + 5y = 13.}$
		\item $\systeme{2x - y = 3, \\  -x + 2y = -5.}$ % \iff \systeme{2x - y = 3, \\  -2x + 4y = -10.}$
		\item $\systeme{4x  - y = -4, \\   x + 5y = 2.}$ % \iff  \systeme{4x  - y = -4, \\   -4x - 20y = -8.}$
		\item $\systeme{-2x - 2y = -5, \\  3x + 7y = 10.}$ % \iff\systeme{-6x - 6y = -15, \\  6x + 14y = 20.}$
		\item $\systeme{7x - y = 10, \\  5x + 2y = 1.}$ % \iff \systeme{35x - 5y = 50, \\  -35x - 14y = -7.}$
	\end{enumerate}
	\end{multicols}
}

\exe{}{
	Pour chaque système de l'exercice \ref{exe:2}, 
		\begin{enumerate}
			\item combiner les équations pour trouver $y$ ; 
			\item trouver $x$ ; 
			\item vérifier que le couple $(x ;y)$ obtenu soit bien solution du système initial.
		\end{enumerate}
}{exe:3}{ \, \\
	\begin{enumerate}[label=Système \roman*) :, itemsep=20pt, leftmargin=80pt]
		\item La somme des deux équations donne
			\[ 8y = 8 \iff y = 1. \]
		En remplaçant $y$ par $1$ dans la première équation, on trouve $x + 3 = 3$, et donc $x =0$.
			\[ (x ; y ) = (0; 1). \]
		
		\item La somme des deux équations donne
			\[ -16y = 16 \iff y = -1. \]
		En remplaçant $y$ par $-1$ dans la première équation, on trouve $x + 7 = 1$, et donc $x =-6$.
			\[ (x ; y ) = (-6; -1). \]
		
		\item La somme des deux équations donne
			\[ 3y = -7 \iff y = -\frac73. \]
		En remplaçant $y$ par $-\frac73$ dans la première équation, on trouve $2x +\frac73 = 3$, et donc $x =\frac13$.
			\[ (x ; y ) = \left(\frac13; -\frac73 \right). \]
		
		\item La somme des deux équations donne
			\[ -21y = -12 \iff y = \frac{12}{21} = \frac47. \]
		En remplaçant $y$ par $ \frac47$ dans la deuxième équation, on trouve $x +\frac{20}7 = 2$, et donc $x =-\frac67$.
			\[ (x ; y ) = \left(-\frac67; \frac47 \right). \]
		
		\item La somme des deux équations donne
			\[ 8y = 5 \iff y = \frac58. \]
		En remplaçant $y$ par $\frac58$ dans la deuxième équation, on trouve $3x +\frac{35}8 = 10 \iff x =\frac{45}{24} = \frac{15}8$.
			\[ (x ; y ) = \left(\frac{15}8; \frac58 \right). \]
		
		\item La somme des deux équations donne
			\[ -19y = 43 \iff y = -\frac{43}{19}. \]
		En remplaçant $y$ par $-\frac{43}{19}$ dans la première équation, on trouve $7x +\frac{43}{19}= 10 \iff x =\frac{147}{133} = \frac{21}{19}$.
			\[ (x ; y ) = \left( \frac{21}{19};  -\frac{43}{19} \right). \]
	\end{enumerate}
}

\exe{}{
	Donner les solutions réelles $(x;y)$ des systèmes suivantes.
	
	\begin{multicols}{2}
	\begin{enumerate}[label=\roman*), itemsep=20pt]
		\item $\systeme{2x+4y=0, \\ 2x -2y = 2.}$
		\item $\systeme{-6x+3y=-3, \\ x -3y = -2.}$
		\item $\systeme{x+y = 2, \\ x+y=2.}$
		\item $\systeme{8x-4y = 6+2x-y, \\ 7x+12y=-5+y.}$
		\item $\systeme{x-y = 1, \\ -x+y=10.}$
		\item $\systeme{\frac12 x - \frac23 y= -1, \\ \frac15 x + \frac72 y= 5.}$
		\item $\systeme{2x - 8y = 2, \\ -4x+16y=-1.}$
		\item $\systeme{5x + y = -2x-1-y, \\ 8x = -1+2y.}$
		\item $\systeme{2y + 12x = -3,\\ -6x=y + \frac32.}$
	\end{enumerate}
	\end{multicols}
}{exe:4}{\, \\
	
	\begin{enumerate}[label=\roman*), itemsep=20pt]
		\item $(x ; y) = (-2 ; 1)$
		\item $(x ; y) = (1;1)$
		\item Les deux équations sont redondantes : on a en fait qu'une seule contrainte pour deux variables et donc un degré liberté. Pour chaque choix de $x$, on peut trouver un $y$ tel que $(x;y)$ soit solution.
		L'ensemble des solutions est
			\[ \{ (x ; y) \text{ tq. } x+y=2, x, y \in \R \} = \{ (x ; y) \text{ tq. } y = -x + 2, \text{ où $x$ parcourt $\R$} \} = \Ccal_f, \]
		où $f(x) = -x+2$.
		\item $(x ; y) = \left(\frac{17}{29} ; -\frac{24}{29}\right)$
		\item Les deux équations sont contradictoires et aucune solution n'existe.
		\item $(x ; y) = \left(-\frac{10}{113} ; \frac{162}{113}\right)$
		\item Les deux équations sont contradictoires et aucune solution n'existe.
		\item $(x ; y) = \left(-\frac{2}{15} ; -\frac{1}{30}\right)$
		\item Les deux équations sont redondantes : on a en fait qu'une seule contrainte pour deux variables et donc un degré liberté. Pour chaque choix de $x$, on peut trouver un $y$ tel que $(x;y)$ soit solution.
		L'ensemble des solutions est
			\[ \left\{ (x ; y) \text{ tq. } -6x -y = \frac32, x, y \in \R \right\} = \left\{ (x ; y) \text{ tq. } y= -6x - \frac32 \text{ où $x$ parcourt $\R$} \right\} = \Ccal_f, \]
		où $f(x) =-6x - \frac32$.
	\end{enumerate}
}


\exe{difficulty=1}{
	Soient $u, v$ deux vecteurs non nuls de $\R^2$ et $\theta$ l'angle dirigé de $u$ vers $v$.
	En adaptant la preuve du théorème \ref{thm:ab-orth}, montrer que
		\[ \scal{u, v} = \norm{u}\norm{v}\cos(\theta). \]
}{exe:prod-scal-2d}{
	todo
}

\exe{difficulty=2}{
	Soient $b_1^*, b_2, b_3 \in\R^3$ trois vecteurs linéairement indépendants.
	\begin{enumerate}
		\item
		Montrer que 
			\[ b_2^* = b_2 - \frac{\japb{b_2, b_1^*}}{\japb{b_1^*, b_1^*}} b_1^* \]
		est orthogonal à $b_1^*$.
		\item
		Montrer que 
			\[ b_3^* = b_3 - \frac{\japb{b_3, b_1^*}}{\japb{b_1^*, b_1^*}} b_1^* - \frac{\japb{b_3, b_2^*}}{\japb{b_2^*, b_2^*}} b_2^* \]
		est orthogonal à $b_1^*$ et $b_2^*$.
	\end{enumerate}
	\emph{Remarque : cet algorithme peut continuer et porte le nom de \emphindex{procédé de Gram-Schmidt}.}
}{exe:gram-schmidt}{
	todo
}


\exe{difficulty=2}{
	Soient $u, v\in\R^n$ deux vecteurs linéairement indépendants.
	\begin{enumerate}
		\item
		Développer la fonction quadratique $f(t) = \norm{u + tv}^2$.
		\item
		Que dire du discriminant du polynôme quadratique $f$ ?
		\item
		En déduire l'\emphindex{inégalité de Cauchy-Schwarz} :
			\[ \abs{\japb{u, v}} \leq \norm{u} \norm{v}. \]
		\item
		Montrer l'égalité suivante pour n'importe quels réels $x_1, \dots, x_n \in \R$.
			\[ |x_1| + \dots + |x_n| \leq \sqrt{n} \sqrt{x_1^2 + \dots + x_n^2} \]
	\end{enumerate}
}{exe:Cauchy-Schwarz}{
	\begin{enumerate}
		\item
		\begin{align*}
			f(t) &= \norm{u + tv}^2, \\
			&= \norm{u}^2 + \norm{tv}^2 + 2\japb{u, tv}, \\
			&= t^2 \norm{v}^2 + 2t \japb{u, v} + \norm{u}^2.
		\end{align*}
		\item
		Le polynôme est toujours positif ou nul car c'est le carré d'une norme.
		Son discriminant est donc négatif ou nul car $f$ n'admet qu'une ou aucune racine.
		\item
		\begin{align*}
			\Delta = b^2 - 4ac &\leq 0, \\
			\left(2\japb{u, v}\right)^2 - 4\norm{v}^2\norm{u}^2 &\leq 0, \\
			4\left(\japb{u, v}^2 - \norm{v}^2\norm{u}^2\right) &\leq 0, \\
			\japb{u, v}^2 &\leq  \norm{v}^2\norm{u}^2, \\
			\abs{\japb{u, v}} &\leq \norm{u} \norm{v}.
		\end{align*}
		\item
		En prenant $u$ le vecteur dont toutes les coordonnées sont 1 et $v$ dont les coordonnées sont $|x_1|, \dots, |x_n|$, l'inégalité de Cauchy-Schwarz donne le résultat attendu.
	\end{enumerate}
}

\exe{difficulty=2}{
	Considérons une série numérique $(a_n)_{n\in\N}$ telle que $\sum_{n\in\N} a_n^2 < \infty$.
	Montrer, à l'aide de l'inégalité de Cauchy-Schwarz (exercice \ref{exe:Cauchy-Schwarz}) que
		\[ \sum_{n\in\N} \frac{|a_n|}{n^s} < \infty \]
	pour tout $s>\frac12$.
}{exe:Cauchy-Schwarz-application}{

}

\shipoutExercise

%\section*{Solutions}
%\addcontentsline{toc}{section}{Solutions}
%
%\shipoutAnswer



